\documentclass[border=10pt]{standalone} 
\usepackage[utf8]{inputenc}
\usepackage{nacal}
\usepackage{pgfplots}
\usepackage{tikz-3dplot}
\usepackage{tikz} 
\usetikzlibrary{matrix,decorations.pathreplacing}
\usetikzlibrary{calc,angles,positioning,intersections,quotes,decorations.markings}
\usepackage{tkz-euclide}
\begin{document}
  \def\Rojo{red!50!black}
  \def\RojoOscuro{red!20!black}
  \def\RojoClaro{red!90!black}
  \def\Verde{green!50!black}
  \def\VerdeOscuro{green!30!black}
  \def\VerdeClaro{green!50!gray}
  \def\Azul{blue!50!black}
  \def\AzulOscuro{blue!35!black}
  \def\AzulClaro{blue!60!gray}

  \begin{tikzpicture}
    [scale=1.05,
     vector/.style={-{Stealth[length=3mm, width=2mm]},very thick,\Azul},]    
    ]
    % vértices del triángulo
    \coordinate (v1) at ( 0 , 0);
    \coordinate (v2) at ( 3 , 0);
    \coordinate (v3) at ( 1 , 2.8);
    \coordinate (v4) at ( 1 , 0  ); % punto intermedio del segmento A
    \coordinate (v5) at (-2 , 2.8); % diferencia de vectores  
    \coordinate (v6) at (0.5, 0 );  % etiqueta X
    
    %\pic [draw,'$\theta$',angle radius=12,angle eccentricity=1.4] {angle = v2--v1--v3};
    \draw pic[-,draw=black,angle radius=15,angle eccentricity=1.4,\Verde,fill=\VerdeClaro!30] {angle = v2--v1--v3};

    \draw[vector] (v1) -- node[below] {A} (v2) node [right]  {$\Vect{a}$};
    \draw[vector] (v1) -- node[left]  {B} (v3) node [above]  {$\Vect{b}$};
    \draw[vector,thick,\AzulClaro] (v1) -- (v5) node [above]  {$\Vect{c}=\Vect{b}-\Vect{a}$};
    % La longitud de C es la longitud de \Vect{b}-\Vect{a}
    \draw[-,\AzulClaro]  (v2) -- node[right] {C} (v3);
    \draw[-,dashed,very thin,gray]  (v3) --  (v5);
    
    % division del triangulo
    \draw[-,gray,dashed,thick]  (v3) -- node[right] {H} (v4);
    % division del segmento A
    \draw[-,dashed,very thick,\Rojo!40]  (v1) -- node[below,\Rojo] {X} (v4);

  \end{tikzpicture}
\end{document}
